\scanpage{171}
\mspage{29}
% slanted
\reversemarginpar\margin{NB $\unsure{\mathscr X|\mathcal X}_{/x}$ inutile}\normalmarginpar

équivalent \deleted{\unsure{encore}} à la suivante e) $\forall\,a$ dans $S$, le \deleted{\unsure{caté}} topos induit $\mathscr X_{/x}$ \deleted{\unsure{est $0$-connexe} et \unsure{irréductible}, \unsure{totalement connexe}} et le topos induit \deleted{\unsure{sur}} $\mathscr X_{/x\times x}$ est \inserted{irréductible.} \deleted{\unsure{e)}} ou encore à la suivante. e') $\forall\,x$ dans $S$, $(S^\wedge_{/x})$ \deleted{\unsure{est irréductible,}} et $S^\wedge_{/x\times x}$ \deleted{est tot. $0$-connexe} sont irréductibles, i.e. $S_{/x}$ et \inserted{$S_{/x\times x}$ sont irréductibles.} \deleted{\unsure{cofiltrants}}

% slanted
\reversemarginpar\margin{$S_{/x}$ inutile, car $\overset{\mathrm{diag}}{x\longrightarrow x\times x}$ est une \unsure{imm.|immersion} ouverte \unsure{non vide}, donc si $x\times x$ est irréd, $x$ l'est aussi. Mais il revient \unsure{donc} au $=$ de dire que les $S_{/x}$ sont irréductibles (condition PS$^{\unsure{x|1}}$1, cf. 2.8, p.~21) \uline{et} que les $x\xrightarrow{\mathrm{diag}}x\times x$ sont dominantes (ce qui implique en effet que $x\times x$ est irréd.) et ceci redonne immédiatement la forme PS$^{x}$2 de la condition cofiltrante.}\normalmarginpar

C'est la forme ``topologique'' la plus faible en apparence que je peux trouver pour la condition $S$ loc cofiltrante, i.e. $\mathscr X$ (\unsure{ou} \deleted{\unsure{Top}} Top$(S)$, ou \unsure{disons|discr.}) \deleted{à la fois loc irréduc et} loc tot. \deleted{\unsure{asph.}} $0$-connexe.

\uline{Définition} 2.13. On dit qu'un topos $\mathscr X$ est \uline{\unsure{localmt|localement}} tot \uline{$0$-connexe} s'il satisfait aux conditions équivalentes a,a',b,d du th. 2.6, i.e. si $\mathscr X$ est loc connexe et si \inserted{de plus} \deleted{\unsure{tout}} topos induit $0$-connexe est tot. $0$-connexe -- \unsure{i.e.} \unsure{encore} le produit fibré de deux objets $0$-connexes sur un troisième est $0$-connexe. Il revient au $=$ de dire que \deleted{\unsure{tz}} \inserted{lorsque \uline{une}} sous catégorie pleine génératrice $S$ de $\mathscr X^\wedge$, formée d'objets $0$-connexes est loc. cofiltrante (en supposant toujours $\mathscr X$ loc \unsure{connexe}, ou ce qui revient au $=$, l'existence d'une $S$ formée d'objets $0$-connexes).

% slanted
\reversemarginpar\margin{C'est confirmé par le fait qu'un topos tot. $0$-connexe est $\unsure{W_\infty|\mathbb W_\infty}$-asph, cf ou $\Rightarrow$ \unsure{ci-bas} prop. 2.3, p. 18 ?}\normalmarginpar

\unsure{§|S} On n'a pas terminé encore la démonstration de 2.6 -- il reste à prouver \inserted{(ce qui \deleted{\unsure{le}} entrainerait)} que si $\mathscr X$ est $0$-connexe, il est $\infty$-connexe -- ou encore que les composantes connexes de $\mathscr X$, supposé loc. tot. $0$-connexe, sont $\infty$-connexes, i.e. \inserted{ou $\infty$-asphériques.$^{*}$}

% archive sheet number: 171
